\chapter{Evaluation}
\label{chap:evaluation}

This chapter reports the empirical evaluation of the Zaguik pipeline. We first describe the experimental setup (Section~\ref{sec:eval-setup}), then present the one-time setup costs (Section~\ref{sec:eval-setup-costs}) and the proof size and timing results across encoding lengths (Section~\ref{sec:eval-results}). The interpretation of these results is deferred to Chapter~\ref{chap:discussion}.

\section{Experimental Setup}
\label{sec:eval-setup}

All experiments were conducted locally on a single machine with a Snapdragon~X 12-core X1E80100 CPU at 3.40\,GHz, 16\,GB of RAM, running Windows~11 (64-bit). The implementation uses Python~3.13.3 and EzKL~22.2.1. No GPU acceleration was used; all proof generation and verification ran entirely on CPU.

\medskip

The structured reference string (SRS) was generated locally using EzKL's facilities rather than retrieved from a remote trusted source. As discussed in Chapters~\ref{chap:system-design} and~\ref{chap:discussion}, this is appropriate for a self-contained research prototype but would differ in a production deployment, where the SRS would be obtained from a public trusted ceremony.

\medskip

Experiments were run sequentially, with no parallel execution across configurations. A single reference output was generated once from a fixed prompt with DistilGPT-2 and reused throughout; it was designed to exceed 2048 bytes so that, for every value of $L$, the encoding window was fully populated (when $L < 2048$, only the first $L$ bytes were encoded and attested). We evaluated five encoding lengths, $L \in \{128, 256, 512, 1024, 2048\}$ bytes, and for each $L$ we repeated proof generation and verification fifteen times ($n = 15$) on this same input, so that the repetitions isolate run-to-run timing variability rather than variation in the data. This yields $5 \times 15 = 75$ proof-generation and verification attempts in total, and the mean values and sample statistics over the fifteen repetitions are reported below.

\medskip

Two clarifications are needed to read these results correctly. First, the pipeline is \emph{gated}: as described in Chapter~\ref{chap:system-design}, the operational regex detector runs in the clear first, and only outputs it judges free of structured PII (emails, phone numbers, identification numbers, and similar) ever reach the proof phase. The attested inputs are therefore ``clean'' in the sense of that detector. Second, and crucially, the circuit policy is stricter and different in kind: it counts the suspicious bytes \texttt{@}, \texttt{-}, \texttt{.}, and the digits \texttt{0}--\texttt{9}, and signals \texttt{PASS} only when this count is zero. Because the period and digits occur in virtually all natural-language text, a multi-hundred-byte LLM output yields a positive count even after passing the regex gate, so the attested result in these runs is a nonzero suspicious-byte count rather than a \texttt{PASS}. This distinction does not affect the quantities measured in this chapter (proof size, proving time, witness time, and verification time are determined by the circuit and the encoding length $L$, not by the value of the count), but it means the experiments should be read as a characterisation of \emph{attestation cost}, not as a measurement of how often outputs satisfy the policy. Accordingly, when we report that all 75 attempts ``succeeded'', we mean precisely that every proof was generated and then verified without error, not that any output produced a \texttt{PASS}.

\section{One-Time Setup Costs}
\label{sec:eval-setup-costs}

Table~\ref{tab:setup-costs} reports the one-time setup costs for the circuit with $L = 2048$, comprising settings generation, circuit compilation, SRS generation, witness setup, and key generation. The total setup time was 102.44\,seconds. This cost is dominated by SRS generation (77.88\,s) and key generation (23.77\,s); the remaining steps each took under one second. Because setup is performed once per circuit and then amortised over all subsequent proofs, it does not contribute to the per-proof cost reported in the next section.

\begin{table}[htbp]
\centering
\caption{One-time setup costs (byte-level PII circuit, $L = 2048$).}
\label{tab:setup-costs}
\vspace{2ex}
\begin{tabular}{l r}
\hline
\hline
Step & Time (s) \\
\hline
Settings generation & 0.14 \\
Circuit compilation & 0.02 \\
SRS generation$^{a}$ & 77.88 \\
Witness (setup) & 0.63 \\
Key generation & 23.77 \\
\hline
Total & 102.44 \\
\hline
\hline
\multicolumn{2}{l}{\footnotesize $^{a}$Generated locally for testing; production deployments would} \\
\multicolumn{2}{l}{\footnotesize obtain the SRS from a trusted ceremony.} \\
\end{tabular}
\end{table}

\section{Proof Size and Timing Results}
\label{sec:eval-results}

Table~\ref{tab:results} presents the proof size, proof generation time, witness computation time, and verification time for each value of $L$, reported as mean $\pm$ sample standard deviation over the fifteen repetitions. All 75 proof-generation and verification attempts completed without error across the five encoding lengths, in the sense clarified in Section~\ref{sec:eval-setup}: every proof was generated and then verified successfully.

\medskip

\begin{table}[htbp]
\centering
\caption{Proof size and timings vs.\ circuit encoding size $L$ (15 independent runs; $n=15$; reference output $\ge 2048$ bytes, prefix-truncated to the $L$-byte window). Each $L$ block shows mean, sample standard deviation ($\pm$Std), and observed range [Min, Max].}
\label{tab:results}
\footnotesize
\setlength{\tabcolsep}{4pt}
\vspace{2ex}
\begin{tabular}{c l r r r r}
\hline
\hline
$L$ (bytes) & & Proof size (bytes) & Proof gen.\ (ms) & Witness (ms) & Verify (ms) \\
\hline
\multirow{3}{*}{128}
  & Mean      & 21896.67 & 15240.07 &  54.31 & 195.64 \\
  & $\pm$Std  & $\pm$27.33 & $\pm$8049.70 & $\pm$21.53 & $\pm$50.65 \\
  & {[Min, Max]} & {[21850, 21944]} & {[10014.94, 33601.85]} & {[40.91, 128.50]} & {[150.74, 305.65]} \\
\hline
\multirow{3}{*}{256}
  & Mean      & 21870.73 & 19240.44 &  89.36 & 200.85 \\
  & $\pm$Std  & $\pm$39.03 & $\pm$9922.94 & $\pm$18.08 & $\pm$65.80 \\
  & {[Min, Max]} & {[21822, 21966]} & {[11019.33, 32606.75]} & {[70.93, 134.03]} & {[149.94, 398.31]} \\
\hline
\multirow{3}{*}{512}
  & Mean      & 21892.67 & 15944.43 & 150.55 & 168.73 \\
  & $\pm$Std  & $\pm$33.82 & $\pm$7903.57 & $\pm$16.24 & $\pm$22.70 \\
  & {[Min, Max]} & {[21835, 21954]} & {[10983.38, 33031.74]} & {[122.49, 183.49]} & {[140.98, 232.17]} \\
\hline
\multirow{3}{*}{1024}
  & Mean      & 21893.07 & 14840.25 & 268.43 & 174.02 \\
  & $\pm$Std  & $\pm$32.81 & $\pm$6573.20 & $\pm$20.57 & $\pm$34.71 \\
  & {[Min, Max]} & {[21841, 21955]} & {[12137.78, 37693.64]} & {[237.38, 295.59]} & {[138.79, 260.96]} \\
\hline
\multirow{3}{*}{2048}
  & Mean      & 21929.27 & 18219.90 & 540.82 & 182.45 \\
  & $\pm$Std  & $\pm$37.22 & $\pm$9359.92 & $\pm$69.33 & $\pm$30.32 \\
  & {[Min, Max]} & {[21867, 21991]} & {[12659.16, 43176.02]} & {[447.87, 704.70]} & {[152.73, 259.90]} \\
\hline
\hline
\end{tabular}
\end{table}

We highlight four factual observations from Table~\ref{tab:results}; their interpretation is taken up in Chapter~\ref{chap:discussion}.

\medskip

\textbf{Proof size is essentially constant.} The mean proof size varies between 21\,871 bytes ($L = 256$) and 21\,929 bytes ($L = 2048$), a spread of less than 0.3\,\% across a sixteen-fold increase in $L$, with no monotonic trend.

\medskip

\textbf{Proof generation dominates and is highly variable.} Mean proof-generation time ranges from roughly 14\,840\,ms ($L = 1024$) to 19\,240\,ms ($L = 256$), with large standard deviations ($\pm$6\,573 to $\pm$9\,923\,ms). It does not vary monotonically with $L$.

\medskip

\textbf{Witness computation grows monotonically with $L$.} Mean witness time increases steadily from 54.31\,ms at $L = 128$ to 540.82\,ms at $L = 2048$, tracking the size of the encoding window.

\medskip

\textbf{Verification time is stable.} Mean verification time stays within a narrow band (roughly 169--201\,ms) across all configurations, with occasional spikes in individual runs but no systematic dependence on $L$.
